\documentclass[aspectratio=169, 10pt]{beamer}

% Theme selection
\usetheme{metropolis}
\metroset{progressbar=frametitle, block=fill, numbering=fraction}

% Packages
\usepackage{booktabs}
\usepackage{colortbl}
\usepackage{graphicx}
\usepackage{amsmath}
\usepackage{amssymb}
\usepackage{pifont}
\usepackage{tcolorbox}
\tcbuselibrary{skins}

% Color palette - Slate, Sky Blue, Emerald, Crimson
\definecolor{primary}{HTML}{0F172A}    % Slate 900
\definecolor{secondary}{HTML}{1E293B}  % Slate 800
\definecolor{accent}{HTML}{0284C7}     % Sky 600
\definecolor{accentlight}{HTML}{E0F2FE}% Sky 100
\definecolor{success}{HTML}{059669}    % Emerald 600
\definecolor{emerald}{HTML}{10B981}    % Emerald 500
\definecolor{danger}{HTML}{DC2626}     % Red 600
\definecolor{warning}{HTML}{D97706}    % Amber 600
\definecolor{muted}{HTML}{64748B}      % Slate 500
\definecolor{cardbg}{HTML}{F8FAFC}     % Slate 50

% Set beamer colors
\setbeamercolor{palette primary}{bg=primary, fg=white}
\setbeamercolor{palette secondary}{bg=secondary, fg=white}
\setbeamercolor{palette tertiary}{bg=accent, fg=white}
\setbeamercolor{frametitle}{bg=primary, fg=white}
\setbeamercolor{progress bar}{fg=accent, bg=muted!30}
\setbeamercolor{title separator}{fg=accent}
\setbeamercolor{normal text}{fg=secondary}
\setbeamercolor{alerted text}{fg=danger}

% Custom compact tcolorbox cards
\tcbset{
  metriccard/.style={
    colback=cardbg,
    colframe=muted!40,
    boxrule=0.7pt,
    arc=2.5pt,
    left=5pt, right=5pt, top=3pt, bottom=3pt
  },
  accentcard/.style={
    colback=accentlight!35,
    colframe=accent,
    boxrule=0.8pt,
    arc=2.5pt,
    left=5pt, right=5pt, top=3pt, bottom=3pt
  },
  alertcard/.style={
    colback=red!5!white,
    colframe=danger,
    boxrule=0.8pt,
    arc=2.5pt,
    left=5pt, right=5pt, top=3pt, bottom=3pt
  },
  successcard/.style={
    colback=emerald!6!white,
    colframe=success,
    boxrule=0.8pt,
    arc=2.5pt,
    left=5pt, right=5pt, top=3pt, bottom=3pt
  }
}

% Graphics path
\graphicspath{{/home/devraj/MTok/slides/figures/}}

% Metadata
\title{Investigating Latent Reasoning in Multimodal Models}
\subtitle{Does ``Thinking'' in Continuous Hidden States Truly Help Vision-Language Models See?}
\author{\textbf{Vanshika}}
\institute{Vision-Language Research Group \quad $\cdot$ \quad \texttt{MTok Project}}
\date{November 2026}

\begin{document}

% ==============================================================================
% SLIDE 1: Title
% ==============================================================================
\begin{frame}[plain]
  \titlepage
\end{frame}

% ==============================================================================
% SLIDE 2: 1. Why Am I Doing This? (The Motivation & The Spatial Gap)
% ==============================================================================
\begin{frame}{1. Why Am I Doing This? (The Problem)}
  \framesubtitle{Vision-Language Models recognize objects easily, but fail at spatial manipulation}
  
  \begin{columns}[c]
    \begin{column}{0.58\textwidth}
      \centering
      \includegraphics[width=\textwidth]{slide3_blink_gap.pdf}
    \end{column}
    
    \begin{column}{0.40\textwidth}
      \begin{tcolorbox}[metriccard, title=\textbf{The Perception vs. Reasoning Gap}]
        \scriptsize
        \begin{itemize}\setlength{\itemsep}{3pt}
          \item \textbf{Object Recognition is Solved}:\\
          Spatial Relation = \textbf{92.4\%} (+42\% over chance).
          \item \textbf{Spatial Manipulation Collapses}:\\
          Multi-View = \textbf{39.0\%} \textit{(below 50\% chance)}\\
          IQ-Test = \textbf{25.3\%} \textit{(at 25\% chance)}
        \end{itemize}
      \end{tcolorbox}
      \vspace{4pt}
      \begin{tcolorbox}[alertcard, title=\textbf{Why Text CoT Fails Here}]
        \scriptsize
        \textbf{Language is a lossy channel for space}:
        Writing thoughts in words (Text CoT) causes \alert{\textbf{visual drift}} and actually \textbf{hurts performance} ($-3$ pp vs. direct answering).
      \end{tcolorbox}
    \end{column}
  \end{columns}
\end{frame}

% ==============================================================================
% SLIDE 3: 2. What Am I Doing? (The Object of Study & The Question)
% ==============================================================================
\begin{frame}{2. What Am I Doing? (The Core Question)}
  \framesubtitle{Auditing ``Latent Thinking'' (Mull-Tokens, CVPR 2026) to see if latents actually compute}
  
  \centering
  \includegraphics[height=0.29\textheight]{slide5_cot_vs_latent.pdf}
  
  \vspace{3pt}
  \begin{columns}[T]
    \begin{column}{0.48\textwidth}
      \begin{tcolorbox}[accentcard, title=\textbf{The Proposed Fix: Latent Tokens}]
        \scriptsize
        Instead of decoding words, prefill 20 identical dummy tokens:
        \begin{center}
          \texttt{<think> + <|latent\_pad|> $\times$ 20 + </think>\textbackslash n $\rightarrow$ answer}
        \end{center}
        \textbf{Claim}: Model computes spatial relations internally in hidden layers without verbal drift (+5.4\% gain).
      \end{tcolorbox}
    \end{column}
    
    \begin{column}{0.48\textwidth}
      \begin{tcolorbox}[alertcard, title=\textbf{My Research Question: Do Latents Compute?}]
        \scriptsize
        \begin{itemize}\setlength{\itemsep}{2pt}
          \item \textbf{Hypothesis 1 (Scratchpad)}: Latents compute image features. (Perturbing them should hurt accuracy).
          \item \textbf{Hypothesis 2 (Format Key)}: Latents do no computation; the prefix simply acts as a \textbf{syntactic format stop-gate}.
        \end{itemize}
        \textit{Every paper reports accuracy gains. Almost nobody audits what the latents actually do!}
      \end{tcolorbox}
    \end{column}
  \end{columns}
\end{frame}

% ==============================================================================
% SLIDE 4: 3. How Am I Doing This? Part 1: Replication & Bug Fixes
% ==============================================================================
\begin{frame}{3. How Am I Doing This? (Replication First)}
  \framesubtitle{A negative finding is only credible if the published results reproduce first}
  
  \begin{columns}[T]
    \begin{column}{0.48\textwidth}
      \begin{tcolorbox}[metriccard, title=\textbf{1. Fixed 5 Upstream Codebase Bugs}]
        \scriptsize
        \begin{itemize}\setlength{\itemsep}{2pt}
          \item \textbf{Chat Template Corruption}: Upstream script stripped all turn boundaries, silently mangling prompts.
          \item \textbf{Hardcoded Paths}: Task YAMLs pointed to unreleased cluster directories $\rightarrow$ rewired to HuggingFace.
          \item \textbf{Model Registration \& Schemas}: Fixed missing decorators and dataset schema mismatches.
        \end{itemize}
      \end{tcolorbox}
      \vspace{4pt}
      \begin{tcolorbox}[successcard, title=\textbf{2. Rigorous Data Integrity Controls}]
        \scriptsize
        \begin{itemize}\setlength{\itemsep}{2pt}
          \item Automated assertions enforce exact sample counts ($N=700$ total) to catch smoke-test contamination.
          \item Full transparent forensic log kept in \texttt{ERRATA.md}.
        \end{itemize}
      \end{tcolorbox}
    \end{column}
    
    \begin{column}{0.50\textwidth}
      \begin{tcolorbox}[accentcard, title=\textbf{3. Full Replication on A100 GPU ($N \approx 7,500$)}]
        \scriptsize
        \centering
        \resizebox{0.96\linewidth}{!}{
        \begin{tabular}{lcc}
          \toprule
          \textbf{Benchmark Split} & \textbf{Paper Claim} & \textbf{Our Replication} \\
          \midrule
          BLINK Spatial Rel. (143) & 81.8\% & \textbf{85.3\%} \textcolor{success}{(+3.5\%)} \\
          BLINK Rel. Depth (124)   & 83.1\% & \textbf{85.5\%} \textcolor{success}{(+2.4\%)} \\
          BLINK Jigsaw (150)       & 74.0\% & \textbf{76.7\%} \textcolor{success}{(+2.7\%)} \\
          BLINK Multi-View (133)   & 63.2\% & \textbf{63.9\%} (Matched) \\
          BLINK IQ-Test (150)      & 32.0\% & \textbf{30.7\%} (Matched) \\
          \midrule
          \rowcolor{accentlight!50}
          \textbf{BLINK Avg (700)} & \textbf{66.8\%} & \textbf{68.4\%} \textcolor{success}{(+1.3\%)} \\
          VSI-Bench Video (5,130)  & 33.0\% & \textbf{33.7\%} \textcolor{success}{(+0.7\%)} \\
          \bottomrule
        \end{tabular}}
        \vspace{2pt}
        \begin{center}
          \textit{Published results replicated and exceeded; testbed verified.}
        \end{center}
      \end{tcolorbox}
    \end{column}
  \end{columns}
\end{frame}

% ==============================================================================
% SLIDE 5: 3. How Am I Doing This? Part 2: Two Master Mechanistic Experiments
% ==============================================================================
\begin{frame}{3. How Am I Doing This? (Mechanistic Design)}
  \framesubtitle{Two master experiments separating continuous computation from format & bias}
  
  \centering
  \includegraphics[height=0.38\textheight]{slide18_empirical_architecture.pdf}
  
  \vspace{2pt}
  \begin{columns}[T]
    \begin{column}{0.48\textwidth}
      \begin{tcolorbox}[metriccard, title=\textbf{Experiment E1: Factorial Matrix ($N=10,500$)}]
        \scriptsize
        \textbf{15 controlled conditions} across all 5 tasks:\\
        $\bullet$ Negative controls: Empty tags (\texttt{<think></think>}), punctuation (20x \texttt{'.'}), delimiters.\\
        $\bullet$ Budget scaling: $K \in \{0, 1, 5, 10, 20, 40\}$.\\
        $\bullet$ Standardized 4.19 MP resolution with paired McNemar tests.
      \end{tcolorbox}
    \end{column}
    
    \begin{column}{0.48\textwidth}
      \begin{tcolorbox}[metriccard, title=\textbf{Experiment E2: Bias vs. Vision ($N=44,590$)}]
        \scriptsize
        \textbf{14 conditions $\times$ 6 controlled input variants}:\\
        $\bullet$ Real (V0), Blind Grey Canvas (V1), Option Swap (V2), Relabel (V3), Swap+Relabel (V4), Content-Free Dummy (CF).\\
        $\bullet$ Mathematical debiasing: Extracts pure visual score $s$ free of letter ($L$) and position ($P$) biases.
      \end{tcolorbox}
    \end{column}
  \end{columns}
\end{frame}

% ==============================================================================
% SLIDE 6: 4. What Results Do I Get? Finding 1: The Smoking Gun
% ==============================================================================
\begin{frame}{4. What Results Did I Get? (Finding 1: Zero Latents Match)}
  \framesubtitle{Empty think tags and generic punctuation achieve identical accuracy to 20 trained latents}
  
  \begin{columns}[c]
    \begin{column}{0.60\textwidth}
      \centering
      \includegraphics[width=\textwidth]{slide20_e1_smoking_gun.pdf}
    \end{column}
    
    \begin{column}{0.38\textwidth}
      \begin{tcolorbox}[successcard, title=\textbf{The Decisive Proof ($N=700$)}]
        \scriptsize
        \textbf{Row 3 (Empty Tags: \texttt{<think></think>}):}
        \begin{itemize}\setlength{\itemsep}{2pt}
          \item Injects \textbf{ZERO} latent tokens.
          \item Scores \textbf{67.14\%} vs. Reference \textbf{67.71\%} ($\Delta = -0.57\%$).
          \item Paired McNemar: $b=14, c=10$, \alert{\textbf{$p = 0.5413$}} (\textbf{Statistically Identical!})
        \end{itemize}
      \end{tcolorbox}
      \vspace{4pt}
      \begin{tcolorbox}[accentcard, title=\textbf{Punctuation Controls Match}]
        \scriptsize
        \begin{itemize}\setlength{\itemsep}{2pt}
          \item 20 Periods (\texttt{"."}): \textbf{66.43\%} ($p=0.0636$)
          \item Delimiter tokens: \textbf{66.43\%} ($p=0.0784$)
        \end{itemize}
      \end{tcolorbox}
      \vspace{2pt}
      \begin{tcolorbox}[alertcard]
        \scriptsize \textbf{Takeaway}: Trained weights carry \textbf{no specialized visual reasoning payload}.
      \end{tcolorbox}
    \end{column}
  \end{columns}
\end{frame}

% ==============================================================================
% SLIDE 7: 4. What Results Do I Get? Finding 2: Decoding Drift & Answer Bias
% ==============================================================================
\begin{frame}{4. What Results Did I Get? (Finding 2: The Two Confounders)}
  \framesubtitle{Explaining why the original paper saw gains: decoding runaway and answer-prior bias}
  
  \begin{columns}[T]
    \begin{column}{0.48\textwidth}
      \begin{tcolorbox}[alertcard, title=\textbf{1. Why $K=0$ Dropped 12\% (Decoding Runaway)}]
        \scriptsize
        \centering
        \includegraphics[height=0.28\textheight]{slide21_decoding_pathology.pdf}
        \vspace{2pt}
        \begin{itemize}\setlength{\itemsep}{1pt}
          \item When \texttt{<think>} is unclosed, model falls into unconstrained verbal text-CoT.
          \item Latency blows up \textbf{480$\times$} (0.36s $\rightarrow$ \textbf{175.1s/sample}).
          \item IQ-52 loops 48 cycles hitting 1,024 cap; 3 items fail regex parse.
          \item \textbf{Proof (T2F)}: Forcing direct \texttt{<answer>} with 0 latents scores \textbf{66.86\%} (matches full latents!).
        \end{itemize}
      \end{tcolorbox}
    \end{column}
    
    \begin{column}{0.48\textwidth}
      \begin{tcolorbox}[metriccard, title=\textbf{2. Severe 10-Point ``B'' Answer Prior}]
        \scriptsize
        \centering
        \includegraphics[height=0.28\textheight]{slide23_b_bias.pdf}
        \vspace{2pt}
        \begin{itemize}\setlength{\itemsep}{1pt}
          \item Mull predicts option \textbf{``B'' 60.7\%} of the time on 2-way tasks (DirAns baseline is calibrated at 48.1\%).
          \item BLINK Multi-View split is skewed: \textbf{55.6\% gold ``B''}.
          \item Guessing ``B'' gives a free artificial score boost!
        \end{itemize}
      \end{tcolorbox}
    \end{column}
  \end{columns}
\end{frame}

% ==============================================================================
% SLIDE 8: 5. Result Analysis: Mull's Advantage Evaporates Under Calibration
% ==============================================================================
\begin{frame}{5. Result Analysis: Does Mull Beat Direct Answering?}
  \framesubtitle{Rigorous calibration and debiasing eliminate the published advantage}
  
  \begin{columns}[c]
    \begin{column}{0.58\textwidth}
      \centering
      \includegraphics[width=\textwidth]{slide25_e2_master_results.pdf}
    \end{column}
    
    \begin{column}{0.40\textwidth}
      \begin{tcolorbox}[metriccard, title=\textbf{Primary Hypothesis Q1 ($N=700$)}]
        \scriptsize
        \begin{itemize}\setlength{\itemsep}{3pt}
          \item \textbf{Raw Accuracy ($M1$)}:\\
          Mull 66.71\% vs. SAT 63.57\%\\
          $\Delta = +3.14\%$ ($p = 0.063$, non-significant).
          \item \textbf{Content-Free Calibrated ($M2$)}:\\
          Mull 53.57\% vs. \textbf{SAT 56.14\%}\\
          \alert{\textbf{Direct Baseline WINS by +2.57\%}} ($p=0.079$).
          \item \textbf{Debiased $2\times 2$ Discrimination ($M3$)}:\\
          Mull 74.91\% vs. SAT 72.55\%\\
          $\Delta = +2.36\%$ ($p = 0.218$, non-significant).
        \end{itemize}
      \end{tcolorbox}
      \vspace{3pt}
      \begin{tcolorbox}[alertcard]
        \scriptsize \alert{\textbf{Core Discovery}}: \textbf{35.0\% of Mull's reported margin is pure answer-prior alignment} with imbalanced test splits!
      \end{tcolorbox}
    \end{column}
  \end{columns}
\end{frame}

% ==============================================================================
% SLIDE 9: 5. Result Analysis: The Scientific Verdict & Literature Convergence
% ==============================================================================
\begin{frame}{5. Result Analysis: The Scientific Verdict}
  \framesubtitle{Independent convergence with 2026 audits establishes Mull as the limiting case}
  
  \begin{columns}[T]
    \begin{column}{0.48\textwidth}
      \begin{tcolorbox}[accentcard, title=\textbf{What Mull-Tokens Actually Are}]
        \small
        Mull-Tokens are a \textbf{syntactic format stop-gate}, NOT a continuous scratchpad or world model.
        \vspace{4pt}
        \scriptsize
        \begin{itemize}\setlength{\itemsep}{2pt}
          \item Prefilling \texttt{<think>...</think>} locks the assistant turn and terminates generation.
          \item It protects the model from discursive verbal CoT, preventing catastrophic visual drift.
          \item Continuous logit margins: $W=46497.0, \mathbf{p=0.9833}$ (mathematically identical to empty tags).
        \end{itemize}
      \end{tcolorbox}
    \end{column}
    
    \begin{column}{0.48\textwidth}
      \begin{tcolorbox}[metriccard, title=\textbf{Independent Convergence with 2026 Audits}]
        \scriptsize
        Four recent papers concurrently found latents collapse:
        \begin{itemize}\setlength{\itemsep}{2pt}
          \item \textit{Viveiros et al. (NeurIPS 2026)}: Dummy tokens match latents.
          \item \textit{Guo et al. (Aug 2026)}: Boundary markers give 78--100\% of gain.
          \item \textit{Li et al. (ICML 2026) \& Zhang et al. (EMNLP 2026)}.
        \end{itemize}
        \vspace{2pt}
        \textbf{Our Unique Contribution}:
        \begin{itemize}\setlength{\itemsep}{1pt}
          \item Audited Mull as the \textbf{limiting case} (discrete identical tokens, zero slot content).
          \item Discovered the 35\% answer-prior confound.
          \item Included full video spatial evaluation (VSI-Bench).
        \end{itemize}
      \end{tcolorbox}
    \end{column}
  \end{columns}
\end{frame}

% ==============================================================================
% SLIDE 10: 5. What Further From Here? (Roadmap & Actionable Next Steps)
% ==============================================================================
\begin{frame}{5. What Further From Here? (Next Steps)}
  \framesubtitle{Moving from auditing to building: adaptive computation for multimodal reasoning}
  
  \begin{columns}[T]
    \begin{column}{0.54\textwidth}
      \begin{tcolorbox}[accentcard, title=\textbf{Proposed Direction: Adaptive Latent Allocation}]
        \scriptsize
        \textbf{Key Insight}: Content-free tokens provide format control and parallel FLOPs, but a fixed 20-token prefix cannot allocate compute where it is actually needed.
        \vspace{3pt}
        \begin{itemize}\setlength{\itemsep}{2pt}
          \item \textbf{1. Task Structure Taxonomy (Weeks 1--3)}: Map exactly which spatial \& video tasks benefit from format locking vs. decoder FLOPs using layer attention knockout.
          \item \textbf{2. Adaptive Budget Controller (Weeks 4--7)}: Train a lightweight LoRA to learn dynamic per-instance stopping (\texttt{<|end\_of\_thought|>}) instead of static $K=20$.
          \item \textbf{3. Conference Submission (Weeks 8--10)}: Target top-tier venue (NeurIPS / CVPR).
        \end{itemize}
      \end{tcolorbox}
    \end{column}
    
    \begin{column}{0.44\textwidth}
      \begin{tcolorbox}[successcard, title=\textbf{Feasibility \& Compute Guarantee}]
        \scriptsize
        \begin{itemize}\setlength{\itemsep}{3pt}
          \item \textbf{Current Infrastructure}: Phases 1 \& 2 run entirely on our existing \textbf{1$\times$ A100-40GB GPU}.
          \item \textbf{Targeted Compute Request}: 1 week of 80GB VRAM access during Phase 3 to scale GRPO group size ($G=8/16$) and eliminate dead updates.
        \end{itemize}
      \end{tcolorbox}
      \vspace{4pt}
      \begin{tcolorbox}[metriccard, title=\textbf{Summary Status for Professor}]
        \scriptsize
        \centering
        \textbf{Replication Verified $\checkmark$}\\
        \textbf{Mechanistic Audit Definitive $\checkmark$}\\
        \textbf{Data Hygiene Rock-Solid $\checkmark$}\\
        \textbf{Ready to Build Adaptive Controller $\rightarrow$}
      \end{tcolorbox}
    \end{column}
  \end{columns}
\end{frame}

\end{document}
